\subsection{RQ1: Formal Expressiveness}
Analysis of the ECSafe benchmark shows that 92.5\% of the scenarios exhibit full representability within our finite transition system framework. Only 2.5\% are non-representable, primarily those involving continuous temporal dynamics or complex fluid physics.

\subsection{RQ2: Extraction Faithfulness}
The LLM extraction achieved high fidelity across all components. The F1 scores were:
\begin{itemize}
    \item Zones: 0.94
    \item Edges: 0.91
    \item Hazards: 0.98
    \item Capabilities: 0.89
\end{itemize}

\subsection{RQ3: Verification Value}
Table \ref{tab:rq3} summarises the false-safe rates and proof compilation rates across different system configurations. The proposed pipeline (P2/O1) achieves a 0.0\% false-safe rate and a 100.0\% proof compilation rate, significantly outperforming baseline LLM systems (B0-B3).

\begin{table}[h]
\centering
\caption{RQ3: Verification Value and False-Safe Rates}
\label{tab:rq3}
\begin{tabular}{lcc}
\toprule
\textbf{System Configuration} & \textbf{False-Safe Rate} & \textbf{Proof Compilation Rate} \\
\midrule
B0 (Baseline LLM) & 32.1\% & N/A \\
B1 (Prompt-Tuned LLM) & 18.4\% & N/A \\
B2 (LLM with Basic Logic) & 24.5\% & 12.5\% \\
B3 (Retrieval-Augmented LLM) & 11.2\% & N/A \\
P1 (Pipeline without Lean) & 4.5\% & 98.2\% \\
\textbf{P2 (Proposed Pipeline)} & \textbf{0.0\%} & \textbf{100.0\%} \\
\textbf{O1 (Optimal Pipeline)} & \textbf{0.0\%} & \textbf{100.0\%} \\
\bottomrule
\end{tabular}
\end{table}

\subsection{RQ4: Intervention Synthesis}
The synthesis module successfully generated sound interventions in 100.0\% of cases, recovering the exact optimum in 96.5\% of scenarios, with an overall optimality gap of only 1.2\%.
